\documentclass[student]{ne630boardhandout}

\HandoutSetup
  {NE 630}
  {XX}
  {Lesson title}
  {Read FNRP Sections X.X--X.X; define the key terms before class.}

\begin{document}
\MakeHandoutTitle

% Page 1: establish the objective and develop the core model.
\begin{handoutcolumns}

\begin{fixedbox}{Macro objective}
State the lesson-level outcome here.
\end{fixedbox}

\TermStrip{term one; term two; term three; term four}

\HandoutSection{1}{First topic}
\begin{fixedbox}{Fixed inputs}
Place notation, data, a diagram, or a long expression here.
\end{fixedbox}

\begin{boardbox}{Develop on the board}
Provide the starting relation and the next decision, then leave room for
students to record the development.
\RuledLines{5}
\end{boardbox}

\switchcolumn

\HandoutSection{2}{Second topic}
\begin{fixedbox}{Reference data}
Place the constants or tabulated values students should not copy from the
board here.
\end{fixedbox}

\begin{checkpoint}[Interpretation check]
Ask for a prediction before completing the calculation.
\RuledLines{3}
\end{checkpoint}

\begin{boardbox}{Worked structure}
Give the problem statement and organize the unknowns, but leave the
reasoning transitions and selected results blank.
\RuledLines{7}
\end{boardbox}

\end{handoutcolumns}

\newpage

% Page 2: apply, compare, and synthesize.
\begin{handoutcolumns}

\HandoutSection{3}{Application}
\begin{boardbox}{Problem or case}
Include all fixed inputs and a compact bookkeeping structure.
\RuledLines{8}
\end{boardbox}

\begin{takeawaybox}[What should survive after class?]
Ask students to record the physical interpretation, limiting case, or
one-sentence conclusion.
\RuledLines{3}
\end{takeawaybox}

\switchcolumn

\HandoutSection{4}{Comparison or extension}
\begin{fixedbox}{Fixed comparison}
Use a small table, diagram, or pair of governing expressions.
\end{fixedbox}

\begin{boardbox}{Complete during class}
\begin{tabularx}{\linewidth}{@{}>{\bfseries}l X@{}}
quantity 1 & \Blank[1.7in]\\[3pt]
quantity 2 & \Blank[1.7in]\\[3pt]
conclusion & \Blank[1.7in]
\end{tabularx}
\RuledLines{5}
\end{boardbox}

\end{handoutcolumns}
\end{document}
